% !PRINTED 356
\weberpage{356}

%<<CONT
\gesperrt{trachtet, so heißt $\mu$ eine Basisform des Ideals $\mathfrak{m}$ in der Ordnung $[Q]$.}

Zum Beweis bemerken wir, daß jede Zahl der Ordnung $[Q]$ wegen (4) in der Form enthalten ist
\weq{10}{\omega = z + \frac{b + \sqrt{D}}{2} y,}
wenn $y$ und $z$ ganze rationale Zahlen sind. Ist nun $\mathfrak{m}$ primär, so ist
\weq{11}{N(\mathfrak{m}) = a}
die kleinste durch $\mathfrak{m}$ teilbare natürliche Zahl, und wenn also $\omega$ durch $\mathfrak{m}$ teilbar sein soll, so muß $z$ durch $a$ teilbar sein. Wenn wir also $z = ax$ setzen, so erhalten wir die Form (9). Umgekehrt ist evident, daß jede Zahl dieser Form durch $\mathfrak{m}$ teilbar ist.

Daß die Darstellung einer Zahl $\mu$ nur auf eine Art in dieser Form möglich ist, folgt daraus, daß $\sqrt{D}$ irrational ist.

Die Form $\mu$ heißt daher eine \gesperrt{Basisform des Ideals $\mathfrak{m}$ in der Ordnung $Q$}. Man erhält daraus eine Basisform desselben Ideals im Körper $\Omega$:
\[ \lambda = ax + \frac{b' + \sqrt{\Delta}}{2} y, \]
wenn $b'$ aus der Kongruenz
\[ b'Q \equiv b \pmod{2a} \]
bestimmt wird, die immer lösbar ist, da $Q$ relativ prim zu $a$, und bei geradem $Q$ auch $b$ gerade ist.

Aus (9) ergibt sich
\weq{12}{N(\mu) = a(ax^2 + bxy + cy^2).}
Die Zahl $\mu$, die man erhält, wenn man in (9) für $x, y$ bestimmte Zahlen setzt, können wir in Idealfaktoren zerlegen:
\[ \mu = \mathfrak{a}\mathfrak{m}, \]
wenn
\weq{13}{N(\mathfrak{a}) = ax^2 + bxy + cy^2}
ist, und es gilt der Satz:

8. Satz. \gesperrt{Ist $\mu$ relativ prim zu $Q$, und haben $x, y$ keine gemeinschaftlichen Teiler, so ist das Ideal $\mathfrak{a}$ primär.}

Wir beweisen ihn so: Es sei $p$ eine in $\mathfrak{a}$ aufgehende rationale Primzahl; dann ist $\mu$ durch $p$ teilbar und daraus folgt zunächst, daß $y$ durch $p$ teilbar sein muß; denn eine Zahl $\omega$ von der Form (10) kann nur dann durch eine in $Q$ nicht aufgehende Primzahl $p$ teilbar sein, wenn sowohl $y$ als $z = ax$ durch $p$ teilbar sind. Demnach kann, da nach Voraussetzung $x$ und $y$
%CONT>>
